% GENERATED FILE. Edit canonical lessons, then run npm run generate:books.
% canonical-source-hash: fnv1a64:fe21b0786ef0a9b7
% canonical-lessons: GE-C494-zanken, GE-C494-schmollen, GE-C494-vertragen, GE-C494-versoehnen, GE-C494-versoehnung

\chapter{Zanken und versöhnen}
\label{ch:ge-zanken-und-versohnen}
\hlchaptermodality{494}

\begin{chapteropening}
\textbf{By the end of this chapter:} \emph{I can talk about quarrels and making up.}

Complete the last lesson of chapter 494: I can talk about quarrels and making up.
\end{chapteropening}

\section[zanken]{zanken — to squabble (sich zanken)}

\label{lesson:GE-C494-zanken}



\begin{quote}
Before the new one: say the German for feeling of safety, then the German for solidarity.
\end{quote}

\subsection*{You'll want to know: zanken}
\textbf{zanken} — \textquotedblleft{}to squabble (sich zanken)\textquotedblright{}.

\begin{quote}
Die Geschwister zanken sich um das letzte Stück Kuchen.
\end{quote}

\begin{quote}
\emph{The brother and sister are squabbling over the last piece of cake.}
\end{quote}

Mostly reflexive: \emph{sich zanken um} + accusative.

\begin{cousinweb}
The origin of \emph{zanken} is not settled; EtymDB lists the German word without a source.
\end{cousinweb}

\begin{grammarlens}[title={What you've built}]
One more word for telling, explaining and describing at B1.
\end{grammarlens}

\subsection*{Your turn}
\begin{itemize}
\raggedright
  \item \emph{Say it:} \emph{zanken}
  \item \emph{Say it:} \emph{zanken}, once more
  \item \emph{Say it:} a sentence of your own with \emph{zanken}
  \item \emph{Recall:} say the German for feeling of safety, then the German for solidarity, then say \emph{zanken} again
\end{itemize}

\begin{tcolorbox}[breakable,colback=teal!4,colframe=teal!35!black,title={Before you move on}]
What does zanken mean? (\textquotedblleft{}To squabble (sich zanken)\textquotedblright{}.) Say it once more.
\end{tcolorbox}
\section[schmollen]{schmollen — to sulk}

\label{lesson:GE-C494-schmollen}



\begin{quote}
Before the new one: say the German for solidarity, then the German for to squabble.
\end{quote}

\subsection*{You'll want to know: schmollen}
\textbf{schmollen} — \textquotedblleft{}to sulk\textquotedblright{}.

\begin{quote}
Nach dem Streit schmollte sie den ganzen Abend.
\end{quote}

\begin{quote}
\emph{After the argument she sulked all evening.}
\end{quote}

Regular: \emph{hat geschmollt}.

\begin{cousinweb}
\emph{schmollen} goes back to Middle High German \emph{smollen}; its older origin is not settled, and EtymDB has no entry for it.
\end{cousinweb}

\begin{grammarlens}[title={What you've built}]
One more word for telling, explaining and describing at B1.
\end{grammarlens}

\subsection*{Your turn}
\begin{itemize}
\raggedright
  \item \emph{Say it:} \emph{schmollen}
  \item \emph{Say it:} \emph{schmollen}, once more
  \item \emph{Say it:} a sentence of your own with \emph{schmollen}
  \item \emph{Recall:} say the German for solidarity, then the German for to squabble, then say \emph{schmollen} again
\end{itemize}

\begin{tcolorbox}[breakable,colback=teal!4,colframe=teal!35!black,title={Before you move on}]
What does schmollen mean? (\textquotedblleft{}To sulk\textquotedblright{}.) Say it once more.
\end{tcolorbox}
\section[vertragen]{vertragen — to get along (sich vertragen); to tolerate}

\label{lesson:GE-C494-vertragen}



\begin{quote}
Before the new one: say the German for to squabble, then the German for to sulk.
\end{quote}

\subsection*{You'll want to know: vertragen}
\textbf{vertragen} — \textquotedblleft{}to get along (sich vertragen); to tolerate\textquotedblright{}.

\begin{quote}
Meine Katze und mein Hund vertragen sich gut.
\end{quote}

\begin{quote}
\emph{My cat and my dog get along well.}
\end{quote}

Reflexive for to get along: \emph{sich vertragen}. Without \emph{sich} it means to tolerate: \emph{Ich vertrage keine Milch.}

\begin{cousinweb}
\emph{vertragen} is inherited from Middle High German \emph{vertragen} and Old High German \emph{fartragan} \textquotedblleft{}to bear, to stand\textquotedblright{} (EtymDB), built on \emph{tragen} \textquotedblleft{}to carry\textquotedblright{}, from Proto-Germanic \emph{draganą}.
\end{cousinweb}

\begin{grammarlens}[title={What you've built}]
One more word for telling, explaining and describing at B1.
\end{grammarlens}

\subsection*{Your turn}
\begin{itemize}
\raggedright
  \item \emph{Say it:} \emph{vertragen}
  \item \emph{Say it:} \emph{vertragen}, once more
  \item \emph{Say it:} a sentence of your own with \emph{vertragen}
  \item \emph{Recall:} say the German for to squabble, then the German for to sulk, then say \emph{vertragen} again
\end{itemize}

\begin{tcolorbox}[breakable,colback=teal!4,colframe=teal!35!black,title={Before you move on}]
What does vertragen mean? (\textquotedblleft{}To get along (sich vertragen); to tolerate\textquotedblright{}.) Say it once more.
\end{tcolorbox}
\section[versöhnen]{versöhnen — to reconcile (sich versöhnen)}

\label{lesson:GE-C494-versoehnen}



\begin{quote}
Before the new one: say the German for to sulk, then the German for to get along.
\end{quote}

\subsection*{You'll want to know: versöhnen}
\textbf{versöhnen} — \textquotedblleft{}to reconcile (sich versöhnen)\textquotedblright{}.

\begin{quote}
Nach einer Woche haben sich die Freundinnen wieder versöhnt.
\end{quote}

\begin{quote}
\emph{After a week the two friends made up again.}
\end{quote}

Reflexive: \emph{sich versöhnen mit} — to make up with.

\begin{cousinweb}
\emph{versöhnen} comes from Middle High German \emph{versüenen} (EtymDB), built on \emph{suone} \textquotedblleft{}atonement, settlement\textquotedblright{}, today \emph{die Sühne}.
\end{cousinweb}

\begin{grammarlens}[title={What you've built}]
One more word for telling, explaining and describing at B1.
\end{grammarlens}

\subsection*{Your turn}
\begin{itemize}
\raggedright
  \item \emph{Say it:} \emph{versöhnen}
  \item \emph{Say it:} \emph{versöhnen}, once more
  \item \emph{Say it:} a sentence of your own with \emph{versöhnen}
  \item \emph{Recall:} say the German for to sulk, then the German for to get along, then say \emph{versöhnen} again
\end{itemize}

\begin{tcolorbox}[breakable,colback=teal!4,colframe=teal!35!black,title={Before you move on}]
What does versöhnen mean? (\textquotedblleft{}To reconcile (sich versöhnen)\textquotedblright{}.) Say it once more.
\end{tcolorbox}
\section[die Versöhnung]{die Versöhnung — reconciliation}

\label{lesson:GE-C494-versoehnung}



\begin{quote}
Before the new one: say the German for to get along, then the German for to reconcile.
\end{quote}

\subsection*{You'll want to know: die Versöhnung}
\textbf{die Versöhnung} — \textquotedblleft{}reconciliation\textquotedblright{}.

\begin{quote}
Zur Versöhnung kaufte er ihr Blumen.
\end{quote}

\begin{quote}
\emph{To make up, he bought her flowers.}
\end{quote}

Plural: \emph{die Versöhnungen}.

\begin{cousinweb}
\emph{Versöhnung} is built on \emph{versöhnen} with \emph{-ung}; the verb comes from Middle High German \emph{versüenen} (EtymDB).
\end{cousinweb}

\begin{grammarlens}[title={What you've built}]
One more word for telling, explaining and describing at B1.
\end{grammarlens}

\subsection*{Your turn}
\begin{itemize}
\raggedright
  \item \emph{Say it:} \emph{die Versöhnung}
  \item \emph{Say it:} \emph{die Versöhnung}, once more
  \item \emph{Say it:} a sentence of your own with \emph{die Versöhnung}
  \item \emph{Recall:} say the German for to get along, then the German for to reconcile, then say \emph{die Versöhnung} again
\end{itemize}

\begin{tcolorbox}[breakable,colback=teal!4,colframe=teal!35!black,title={Before you move on}]
What does die Versöhnung mean? (\textquotedblleft{}Reconciliation\textquotedblright{}.) Say it once more.
\end{tcolorbox}
